\chapter{The Genesis Protocol: The Birth of a Node}

The preceding chapters described the Sovereign Stack as it looks when it is mature: seven layers, each doing its part, each speaking only to its neighbours. No community begins there. A Node is not switched on. It is grown, slowly, by people who first have to learn to trust one another, then find a place, then build, and only then teach their tools to help.

It helps to be honest about the odds. Diana Leafe Christian, who spent years studying people who tried to found intentional communities, estimated that probably nine in ten such groups never get off the ground \parencite{christian2003}. They founder on money, on land, on hidden expectations, and above all on the difficulty of people deciding things together. The Genesis Protocol exists to improve those odds. It does not guarantee success, and no protocol could. It is a suggested path, drawn from what has worked and failed for others, that a founding group can follow, adapt, or set aside.

The protocol runs from the human outward. People come first, then their agreement, then a legal shield, then a place, then the first hardware, and only then the cybernetic layers that make a Node more than a group of neighbours. Earlier drafts of this chapter placed servers first, on the theory that a community needed silicon before it could think. That ordering is withdrawn. Communities think with their people. The Stack arrives to help them remember, measure, and coordinate once there is something worth coordinating.

This chapter also faces something the earlier drafts left out. Communities that hold land, feed themselves, and run their own infrastructure have often been met with hostility, sometimes with violence. A protocol for genesis that ignored that history would be naive. This one plans for it, and it answers it with strategic nonviolence and a set of instruments of protection, chosen because they work and because they keep a young community safe, legitimate, and whole.

\section{Principles of Genesis}

The following principles run through every phase.
\begin{itemize}
  \item \textit{People before platforms.} Trust, agreements, and habits of working together come before hardware or legal entities.
  \item \textit{Start where you are.} A genesis can begin this weekend, with no servers, no money, and no land.
  \item \textit{Covenant, not calculation.} A community is formed by people choosing one another and making promises, not by an algorithm matching profiles.
  \item \textit{Honest from the first day.} With neighbours, with authorities, with game players, and with one another.
  \item \textit{Strategic nonviolence.} The community acts without violence because nonviolence is more effective, safer for its most vulnerable members, and more sustainable.
  \item \textit{Protection without militarisation.} The community prepares to protect its people, and it guards against its protectors becoming a police force.
  \item \textit{Grow at the pace of trust.} Each phase begins when the group is ready for it, not when a schedule says so.
  \item \textit{Design for an honourable ending.} Every genesis provides, from the start, for how it will close if it must.
\end{itemize}

\section{Three Ways a Genesis Begins}

A genesis can start in three different settings. The phases that follow are the same in each, but their weight differs.

\subsection{In Oasis}

In the Oasis simulation, the Genesis Protocol is the game's main progression arc. Players found a settlement, write a compact, build, and govern, all with simulated people, land, and hardware. Nothing they do touches the real world. It is practice, and for many people it will be the first time they have seen what founding a community involves.

\subsection{From the Legacy World}

A group of neighbours, friends, or coworkers living in the legacy economy decides to found a Node. This is the setting the rest of this chapter follows most closely, because it is the hardest: the group starts with nothing but itself.

\subsection{From an Existing Node}

An established Node, or several, may help found a new one nearby, for example when a neighbouring village asks for help or a Node grows too large for its councils. This begins as a proposal ratified in the existing Node's Layer 5, under the federation and nesting arrangements of Chapter 6. The established Node can lend skills, equipment, and experience. It cannot found the new Node on its behalf. The new community still gathers, writes and accepts its own Genesis Compact, and governs itself from its first day.

\section{The Phases at a Glance}

Table~\ref{tab:genesis-phases} summarises the phases. The durations are indicative only: some groups move faster, many move slower, and a group that rushes a phase usually pays for it later.

\begin{table}[htbp]
\centering
\small
\begin{tabular}{|p{0.07\textwidth}|p{0.17\textwidth}|p{0.42\textwidth}|p{0.20\textwidth}|}
\hline
\textit{Phase} & \textit{Name} & \textit{What happens} & \textit{Indicative duration} \\
\hline
0 & Gathering & Shared meals, mutual aid, a garden or a tool library; learning how the group decides and disagrees. & Months \\
\hline
1 & The Compact & Founders write and unanimously accept the Genesis Compact and mark it together. & Months \\
\hline
2 & The Shield & A lawful cooperative entity is formed; protection and legal readiness begin. & Months \\
\hline
3 & A Place & Land or space is secured, in right relationship with its history and its neighbours. & Months to years \\
\hline
4 & The Anchor & People install the first meters, sensors, and modest computing; Layers 1 to 3 come alive. & Months \\
\hline
5 & Ignition & The first workflows are drafted, ratified, and run, with autonomy granted in stages. & Months \\
\hline
6 & Maturation & New members, more councils, founder roles sunset, federation with other Nodes. & Years, ongoing \\
\hline
\end{tabular}
\caption{The phases of a genesis.}
\label{tab:genesis-phases}
\end{table}

\section{Phase 0: Gathering}

The first phase needs no technology at all. A group that wants to found a Node starts by doing, together, the kinds of things a Node will eventually do: sharing meals, starting a community garden, opening a tool library or a repair café, organising mutual aid for neighbours in need. Movements such as Food Not Bombs, which has shared free food since 1980, and Occupy Sandy, which organised mutual aid after Hurricane Sandy in 2012, show how much can be done this way with no formal structure at all. A shared notebook can serve as the first ledger, and a weekly meeting as the first council.

This phase matters less for what it produces than for what it reveals. A group discovers who shows up, how decisions actually get made, what happens when two people want incompatible things, and whether people still like one another after a hard week. Jo Freeman warned that groups without explicit structure do not escape hierarchy; they simply hide it, so that power flows to whoever is best connected and least accountable \parencite{freeman1972}. Even in Phase 0, then, a group should write down how it decides and how it handles disagreement, and revisit those notes as it learns.

Phase 0 is open to everyone. A Node does not require its founders to own land or hold capital, and some of the most important Nodes may be founded by renters, by people without stable housing, or by migrants whose legal status makes formal ownership hard. Nothing in this phase asks anyone for money or for papers. Oasis can help here too, as a sandbox in which a group can practise governance and see its dynamics before anything real is at stake. But the essential work is face to face.

Celebrate as you go. A harvest meal, a repair-café anniversary, a song written for the garden: these are not distractions from genesis. They are how a group becomes a people, and how it remembers later why it began.

\section{Phase 1: The Genesis Compact}

When a group is ready to commit, it writes its Genesis Compact: the founding agreement described in Chapter 6, which defines how the community governs itself and what it will never do. Rosabeth Moss Kanter's study of nineteenth-century American communes found that the communities that lasted were those that built strong commitment through shared investment, shared ritual, and a clear sense of who belonged \parencite{kanter1972}. The Compact is where that commitment first takes form.

\subsection{Writing the Compact}

The Compact is written by the founders themselves, in plain language, over as many sessions as it takes. It contains the entrenched core that every Node shares (Chapter 6): the unconditional survival baseline, the existence of ecological floors, the right to refuse work, the right of exit, the privacy of ballots, self-reports, and care leave, due process, and the non-adversarial treatment of Non-Stewards. Table~\ref{tab:genesis-compact} lists the further commitments this chapter recommends for a founding group.

\begin{table}[htbp]
\centering
\small
\begin{tabular}{|p{0.27\textwidth}|p{0.63\textwidth}|}
\hline
\textit{Commitment} & \textit{Purpose} \\
\hline
Nonviolent discipline & The community's collective actions are nonviolent; members are trained in de-escalation; breaches are addressed through the restorative process. \\
\hline
Protection charter & Defines what community safety roles may and may not do, to whom they answer, and how they rotate (Section~\ref{sec:genesis-protection}). \\
\hline
Safeguarding & Rules protecting children and vulnerable members, including an external contact for complaints that bypasses founders. \\
\hline
Founder sunset & Founding roles and any special powers expire on a fixed date; no founder holds a veto. \\
\hline
Money and exit & What each founder contributes, how it is recorded, and how it is returned if they leave or the Node dissolves. \\
\hline
Land and neighbours & Commitments to right relationship with the land's history and to good faith toward neighbours (Phase 3). \\
\hline
Dissolution & How the community would close if it had to, and where its assets would go (Section~\ref{sec:genesis-ending}). \\
\hline
\end{tabular}
\caption{Recommended additions to the Genesis Compact for a founding group.}
\label{tab:genesis-compact}
\end{table}

\subsection{Accepting the Compact}

Each founder accepts the Compact, and acceptance must be unanimous: nobody is bound by a constitution they did not agree to. People who join later accept it as a condition of membership, and the right of exit keeps that acceptance meaningful.

The acceptance should be marked. How it is marked belongs to the community: a meal, a reading under a tree, a song, a planting, a ceremony drawn from the founders' own traditions. What matters is that everyone is present, that the promises are said aloud, and that the moment is remembered. The Compact is a covenant between people before it is a record in a ledger. When Layer 5 comes online in Phase 4, the Compact, first kept on paper, is entered as the community's root record, and the founders sign it again with their keys.

\subsection{Steward Genesis}

Each founder also sets up their identity in the Stack. In Oasis this step is called Character Creation, and earlier drafts carried that name into the real Node, along with an interview in which a language model mapped each founder's ``psychological archetypes'' and ``metabolic needs'' and then ``mathematically aligned'' the group. That design is withdrawn. It would build exactly the profile of each person that Chapter 6 forbids, and it mistakes a community for a team assembled by matching.

In a real Node, Steward Genesis is private and in each person's own hands. Each founder generates their keys on their own device. They choose a few trusted people to hold shares of a recovery secret, confirmed in person, as Chapter 6 describes. If they wish, they may record skills they are willing to offer, as self-attestations that others can later confirm by watching them work. Nothing about their psychology, health, or body is collected. Their needs are met through the unconditional survival baseline, which asks nothing of them.

No one begins with Reputational Wealth. It is earned by delivering value to the community, and at founding no one has yet delivered value within the Node. Every founder starts with equal membership standing and a vote in the commons chamber. Until Reputational Wealth has accumulated, the Compact provides that decisions are made in the commons chamber alone.

\subsection{Founders, Power, and Harm}

Young communities are especially vulnerable to concentrated power. Founders often hold informal influence far beyond their formal role: they know the history, they made the first contacts, they may have put in the most money. Charismatic founders can come to dominate a group's thinking, and the same intensity that makes a founding group feel like a family can make it hard to name abuse when it happens. Decaying Reputational Wealth helps, but it does not reach informal power. The Compact therefore includes explicit protections: founder roles sunset on a fixed date; no founder holds a veto; facilitation and key roles rotate; safeguarding rules protect children and vulnerable members; and there is an external person or organisation to whom anyone can raise a concern without going through the founders. A community that cannot be criticised by its own members is already in danger.

\section{Phase 2: The Shield}

Only after the Compact exists does the group form its legal entity. Forming a company among people who have not yet agreed how they will decide things traps them in obligations before they know how they handle conflict.

The entity is the lawful, cooperative-style body described in Chapter 8: a worker or multi-stakeholder cooperative, a Social Purpose Corporation, a community land trust, or a combination. Its governing documents mirror the Compact: member control, an asset lock, officers bound to ratified decisions, and a dissolution clause consistent with the Compact's. Earlier drafts described this entity as a ``Trojan Horse''; that framing is retired. The entity is what it says it is.

Founders' contributions are recorded and returned on the terms the Compact sets. If the Node needs fiat early on, it raises it under the rules of Chapter 8: only against a ratified needs list, from genuine surplus, with honest dealing and stop rules. Internal credentials are never ``translated'' into legal licences. Where a role requires a licence, the entity pays for a member's real training and examination.

The Shield is also where protection begins. Early tasks for the new entity include opening a legal defence fund, building a relationship with a lawyer who understands land, housing, and civil rights, and making contact with local legal-observer and solidarity networks.

\section{Phase 3: A Place}

A Node is a place. Securing one is often the longest and most expensive phase of a genesis.

\subsection{Paths to Land and Space}

There is more than one way to hold a place, and Table~\ref{tab:genesis-land} compares the common paths.

\begin{table}[htbp]
\centering
\small
\begin{tabular}{|p{0.22\textwidth}|p{0.34\textwidth}|p{0.34\textwidth}|}
\hline
\textit{Path} & \textit{Strengths} & \textit{Cautions} \\
\hline
Purchase by the cooperative entity & Secure tenure; full control of land use. & Requires capital; favours groups that have it. \\
\hline
Community land trust & Removes land from speculation permanently; long precedent, from New Communities in Georgia (1969) to the Dudley Street Neighborhood Initiative in Boston. & Needs a trust partner or the capacity to found one. \\
\hline
Lease or agreement with an owner or municipality & Low capital; can use idle land or buildings; can be a bridge to ownership. & Tenure may be insecure; terms must protect the community's investment. \\
\hline
Converting an existing neighbourhood & Renters' associations, housing cooperatives, and neighbours on a street can found a Node where they already live. & Requires broad neighbour participation and patience with existing ownership patterns. \\
\hline
Occupation of unused land or buildings & Has given landless people a foothold where no other path existed, as Brazil's Landless Workers' Movement has shown since 1984. & Unlawful in most jurisdictions; exposes members to eviction, arrest, and violence; must be a deliberate collective decision. \\
\hline
\end{tabular}
\caption{Paths to land and space.}
\label{tab:genesis-land}
\end{table}

This protocol's default is lawful tenure, because it gives a young community the security to invest in soil and infrastructure and spares its members avoidable risk. It does not condemn occupation. Occupation of unused land is a long-standing form of nonviolent direct action, and in some places it has led to lasting, legal settlements. A community that considers it should treat it as what it is: a high-risk collective act. It should be decided under the Compact with full knowledge of the legal and physical risks, carried out within the community's nonviolent discipline, and never imposed on members, such as those with precarious legal status or responsibility for children, who cannot safely bear those risks.

\subsection{Right Relationship with Land}

No land is empty. Before acquiring it, a founding group should learn its history: who lived there, who lives there still, who was removed, and who holds rights or relationships to it that legacy title does not record. Where Indigenous nations have ties to the land, the group should seek them out and listen, ready to honour access rights, co-stewardship arrangements, or a return of land if that is what right relationship requires. A statement acknowledging the land's history is a beginning, never a substitute for these relationships.

The land deserves listening too. Before any sensor is installed, the founders should spend time, ideally a full cycle of seasons, observing where water flows, where frost settles, what grows, and what lives there. The ecological floors the community will later ratify should rest on that patient observation as much as on any model.

\subsection{Codes, Permits, and Interconnection}

Zoning rules, building codes, septic and water permits, and the conditions utilities impose on connecting generators to the grid all apply to a new Node. They take time, often far more than founders expect. They are also, mostly, written in response to real dangers: fires, collapses, contaminated wells. A Node that builds to code protects its people and denies hostile officials an easy pretext. Where rules obstruct sound ecological building, the community can seek variances, work with sympathetic officials, and advocate for change, but it should begin by understanding the rules well.

\section{Phase 4: The Thermodynamic Anchor}

With a place secured, people begin installing the first physical infrastructure: meters on the water supply and the electrical panel, sensors in soil and storage, perhaps a few solar panels and a battery. This is when the Stack first comes alive.

The first hardware is modest. Salvaged computers and microcontrollers are enough to run the early layers. Layer 2 devices are joined to a local mesh that never touches the legacy internet, and each device is registered with its own hardware identity. Layer 3 begins recording signed measurements. When it can verify that useful physical work has been done, such as water purified or energy stored, it mints the Node's first Exergy Credits under Proof-of-Thermodynamic-Work, to the commons, subject to demurrage from the first moment. Earlier drafts had the Orchestrator mint these first credits; that is not its role. Layer 3 mints. The Orchestrator never does.

In this phase nothing moves on its own. Valves and switches stay under manual control. Layer 5 comes online to hold the Compact as the community's root record, together with each founder's keys. Language models and other heavier tools, if the community wants them, can come later; the Stack's core functions work without them.

\section{Phase 5: Ignition}

Ignition is the moment the Orchestrator begins to act, and it happens only when the community asks it to.

The founders, using Layer 6's draft, read-back, and sign process, write their first workflows as BPMN process diagrams and ratify their first standing instruments in Layer 5: the ecological floors, the unconditional survival baseline, the priority order, the human capacity ceiling, and the first emergency playbooks. Before ratifying, they can rehearse each workflow in Oasis's shadow simulator, whose telemetry is tagged as simulated so that it can never be mistaken for the real thing. Only then does the Orchestrator run the ratified workflows, within the actuation envelopes they define.

Autonomy is granted in stages. A new workflow first runs in observation mode, reporting what it would do without doing it. When the community trusts its judgement, it moves to advisory mode, proposing actions for a person to approve. Only then does it receive bounded authority to act on its own within its envelope. Each step is a ratified decision, and the community can step back at any time.

\section{Phase 6: Maturation}

As the Node becomes a good place to live, others ask to join. Each new member accepts the Compact, sets up their identity as the founders did, and receives equal membership standing. Membership boundaries and the process for joining are set by the community's councils, as Chapter 6 describes.

The single founding council gives way to many. Water, energy, food, care, and other domains each gain councils of the people who depend on them, nested and overlapping as Ostrom's work suggests they should be. Founder roles sunset as the Compact provides. Reputational Wealth earned in the early years decays, so that newer members' contributions carry their proper weight. Future Generations panels begin reviewing long-horizon decisions. The Node joins the federation of other Nodes, and its Layer 7 begins measuring, through the Severance Protocol of Chapter 8, how its dependence on the legacy economy declines.

Earlier drafts said that at this stage a Node is ``fully autonomous.'' That is not how it should be described. A mature Node is increasingly self-reliant in energy, water, food, and care, and increasingly able to govern itself well. It still trades honestly with its neighbours, still depends on the wider mesh, and still lives within a legacy state. Autonomy is a direction of travel, not a finish line.

\section{Meeting Legacy Aggression}
\label{sec:genesis-aggression}

The legacy world will not always welcome a Node. The more a community holds land, feeds itself, and runs its own systems, the more likely it is to attract hostility from officials, from powerful neighbours, and occasionally from the state's armed forces.

\subsection{Forms of Violence}

Johan Galtung distinguished \textit{direct} violence, harm done by identifiable people to identifiable people, from \textit{structural} violence, harm built into institutions and the distribution of power, which can damage lives without anyone striking a blow \parencite{galtung1969}. A new Node may meet both. Direct violence includes evictions, raids, harassment, arson, and attacks by hostile individuals or groups. Structural violence is more common and often more effective: selective code enforcement, zoning refusals, denial of utility connections, the closing of bank accounts, and the gradual criminalisation of practices the community depends on. A third form is covert: surveillance, infiltration, and disinformation designed to divide a community from within or discredit it outside.

\subsection{What History Teaches}

None of this is hypothetical. Table~\ref{tab:genesis-cases} lists documented cases and the lessons they offer a founding group.

\begin{table}[htbp]
\centering
\small
\begin{tabular}{|p{0.22\textwidth}|p{0.36\textwidth}|p{0.32\textwidth}|}
\hline
\textit{Case} & \textit{What happened} & \textit{Lesson} \\
\hline
COINTELPRO, United States (1956--1971) & FBI programme of surveillance, infiltration, and disinformation against political and civil-rights movements, exposed in 1971 and investigated by the Senate's Church Committee. & Expect infiltration and disinformation; minimise what can be surveilled; resist turning suspicion into internal witch-hunts. \\
\hline
MOVE, Philadelphia (1985) & Police dropped a bomb on a communal household during an armed standoff; the fire killed eleven people, five of them children, and destroyed sixty-one homes. & State violence can be catastrophic, and children suffer first; isolation from neighbours removed every brake on escalation. \\
\hline
Christiania, Copenhagen (since 1971) & A self-governing community in former barracks endured decades of raids and legal pressure, then negotiated a collective purchase of the land in 2012. & Persistence, public support, and negotiation can turn precarious tenure into lasting security. \\
\hline
MST, Brazil (since 1984) & The Landless Workers' Movement has settled hundreds of thousands of families on unused land; in 1996 military police killed nineteen of its members at Eldorado dos Carajás. & Occupation can win land, but carries mortal risk; legal grounding and broad alliances matter. \\
\hline
Standing Rock, United States (2016--2017) & Opposition to the Dakota Access Pipeline met water cannon in freezing temperatures and surveillance by private security contractors. & Expect private as well as public force; documentation and legal observation are essential. \\
\hline
ZAD, Notre-Dame-des-Landes (2018) & After the state abandoned a planned airport, thousands of gendarmes evicted parts of the occupation; some projects were later regularised through agreements. & Even victory can bring eviction; prepared negotiators can secure a legal future for part of a project. \\
\hline
\end{tabular}
\caption{Documented cases of legacy aggression and their lessons.}
\label{tab:genesis-cases}
\end{table}

\subsection{Levels of Pressure}

It helps to name the level of pressure a Node is under, so that its response can match it. Table~\ref{tab:genesis-pressure} uses the four regimes the Oasis simulation models.

\begin{table}[htbp]
\centering
\small
\begin{tabular}{|p{0.17\textwidth}|p{0.33\textwidth}|p{0.40\textwidth}|}
\hline
\textit{Level} & \textit{Typical signs} & \textit{Suggested posture} \\
\hline
Below the radar & Little official attention; curious neighbours. & Build neighbour relationships; get codes and paperwork right; open the legal defence fund; train in de-escalation. \\
\hline
Under scrutiny & Inspections, complaints, media interest, unexplained visitors. & Open days and transparency; documentation; legal advice on every notice; prepare evacuation and continuity plans. \\
\hline
Adversarial & Enforcement actions, account closures, harassment, threats. & Activate solidarity networks and legal observers; public communication; rehearse plans; move vulnerable members to safety if needed. \\
\hline
Active siege & Eviction operation, raid, or sustained attack. & Nonviolent discipline; children and vulnerable people away first; documentation and legal response; no physical resistance that endangers lives; negotiate. \\
\hline
\end{tabular}
\caption{Levels of legacy pressure and suggested postures.}
\label{tab:genesis-pressure}
\end{table}

\section{Strategic Nonviolence}

\subsection{Why Nonviolence}

The Collective chooses nonviolence for strategic reasons. Gene Sharp argued that all political power rests on the cooperation of many people, and that organised non-cooperation can withdraw it; he catalogued 198 methods of nonviolent action, from protest and persuasion to strikes, boycotts, and the building of parallel institutions \parencite{sharp1973}. Erica Chenoweth and Maria Stephan studied 323 major resistance campaigns between 1900 and 2006 and found that nonviolent campaigns succeeded about twice as often as violent ones, largely because they drew far broader participation and were more likely to win over defectors from the other side \parencite{chenoweth2011}.

For a Node, the strategic case is sharper still. A community with children, elders, and gardens cannot win an armed confrontation with a state and should never try. Its strongest protection is legitimacy: neighbours who know it, officials who have no pretext, journalists who can see that it is peaceful, and courts that find it lawful. Violence by the community would destroy all of these at once. Nonviolence protects the people most at risk, keeps the community united, and preserves the conditions under which it can keep growing.

\subsection{Nonviolence Is Not Passivity}

Strategic nonviolence is active. It includes bold, visible, joyful action: festivals and open days that invite the neighbourhood in, public celebrations of a first harvest, solidarity with neighbours facing eviction, mass non-cooperation with unjust orders, and civil disobedience accepted openly with its consequences. Most of all, it includes the constructive work of building the alternative itself. The Node is its own most powerful nonviolent action: every meal it serves and every kilowatt it stores is proof that another way of living is possible.

\subsection{Nonviolent Discipline}

Nonviolence fails when it is improvised. The Compact commits the community to nonviolent discipline in all its collective actions, and members prepare for that commitment as they would for any other skill: de-escalation training, role-play of tense encounters, clear agreements about who speaks for the community in a confrontation, and rehearsed plans for evictions and raids. When someone breaches the discipline, the community responds through the restorative process of Chapter 6, and the person may be asked to step back from protective or public roles.

\section{Instruments of Protection}
\label{sec:genesis-protection}

Nonviolence does not mean being unprepared. Table~\ref{tab:genesis-instruments} lists the instruments of protection this protocol recommends, the phase in which each should begin, and the risk to watch in each.

\begin{table}[htbp]
\centering
\small
\begin{tabular}{|p{0.22\textwidth}|p{0.34\textwidth}|p{0.08\textwidth}|p{0.24\textwidth}|}
\hline
\textit{Instrument} & \textit{What it does} & \textit{Phase} & \textit{Watch for} \\
\hline
Neighbour alliances & Relationships, shared projects, and mutual aid with the surrounding community; the strongest deterrent. & 0 & Treating neighbours as a shield rather than as people. \\
\hline
Transparency and open days & Lets neighbours, officials, and media see what the Node is and does. & 2 & Exposing members' private lives. \\
\hline
Legal defence fund and counsel & Money and a lawyer ready before a crisis. & 2 & Reliance on one lawyer or one funder. \\
\hline
Know-your-rights training & Members know their rights in inspections, raids, and evictions. & 2 & Assuming rights will be respected. \\
\hline
Legal observers and documentation & Trained observers record events; footage and notes support legal defence. & 2--3 & Recording laws vary; footage can expose members. \\
\hline
Community safety team & Unarmed, trained, rotating members who de-escalate, keep watch, and coordinate in emergencies. & 3 & Drift toward policing members. \\
\hline
Unarmed civilian protection & Partners who accompany communities under threat, following the model of organisations such as Nonviolent Peaceforce. & 3 & Dependence on outsiders' presence. \\
\hline
Evacuation, continuity, and child protection & Rehearsed plans to move children and vulnerable people to safety and keep essential functions running elsewhere. & 3 & Plans that exist only on paper. \\
\hline
Physical safety & Fire safety, sound structures, code compliance. & 3--4 & Treating safety as optional because officials are hostile. \\
\hline
Digital and operational security & Data minimisation, local systems, no central registries of members or locations. & 4 & Secrecy that corrodes openness. \\
\hline
Safeguarding & Protection of children and vulnerable members from harm inside the community. & 1 & Founders investigating complaints against themselves. \\
\hline
\end{tabular}
\caption{Instruments of protection.}
\label{tab:genesis-instruments}
\end{table}

\subsection{Nuances}

Every instrument of protection carries a risk of becoming a danger in its own right. The following nuances should shape how a community uses them.

\subsubsection{Protection Without Militarisation}

A community safety team can slowly become a police force: first a radio and a schedule, then uniforms, then the authority to tell other members what to do. To prevent that drift, the protection charter in the Compact defines the team's mandate narrowly. Its members hold no powers beyond those any member holds. They serve short, rotating terms, carry no weapons, and answer to a council rather than to themselves. Any change to the charter is a Constitutional-tier decision.

\subsubsection{Who Watches the Protectors}

Complaints about protective roles go to a body independent of them, such as the Review Council of Chapter 6, whose members are drawn by lot. Protective roles are reviewed regularly, and anyone who has served in one can be asked to step down.

\subsubsection{Security Culture Without Paranoia}

The best security is having little to hide and little to steal. Data minimisation, the pairwise identifiers of Chapter 6, and the records boundary of Chapter 8 protect members far better than suspicion. A community that treats every newcomer as a spy will lose the openness that makes it worth defending.

\subsubsection{Infiltration and False Accusation}

Infiltrators exist, and so does the damage done by unfounded accusations. During COINTELPRO, casting false suspicion on genuine members was itself a tactic, sometimes called snitch-jacketing. A community therefore responds to behaviour, not to suspicion. Concerns go through due process, never public denunciation. Someone who repeatedly pushes the group toward violence or illegality, or who breaches safeguarding rules, is addressed for what they did, whoever they turn out to be.

\subsubsection{Protecting People and the Nonviolent Stance}

The Compact's nonviolent discipline governs what the community does together. It does not and cannot override the rights the law gives each individual, including the right to defend oneself or another person from immediate harm. The community does not organise armed defence, and it does not stockpile weapons. Whether and how individual members may keep lawful firearms is a question of law for each person. The community may still decide, as a Compact matter, to keep certain spaces weapon-free, such as assemblies, childcare spaces, and any collective action. During any confrontation, the first duty is to the most vulnerable. Children and people who cannot safely face risk are moved away first and are never placed in front of a confrontation.

\subsubsection{Individual Rights and the Collective Posture}

Participation in any risky action is voluntary. Members whose legal status is precarious, who care for dependants, or who are especially vulnerable to state violence are never pressed into high-risk roles, and their refusal carries no stigma. The community's posture is collective, but the risks fall on individuals, and the community must honour that.

\section{When a Genesis Ends}
\label{sec:genesis-ending}

Most attempts to found a community do not succeed, and some Nodes that do succeed will one day close. An ending is not a failure of character, and the Compact prepares for it from the start.

When a community decides to dissolve, under the procedure its Compact sets, it honours its obligations to creditors, members, and neighbours. It returns members' contributions on the agreed terms. Under the asset lock, it passes its remaining assets, including land, tools, and equipment, to another commons, a neighbouring Node, or a community land trust, never to individuals. It revokes keys and deletes personal records, except those the law requires the legal entity to keep. It hands over care of the land to someone who will continue it. And it records what it learned, in a commons of lessons that other founding groups can read, so that its experience strengthens the next attempt.

An ending also deserves to be marked. People who gave years to a shared dream will grieve it. A closing gathering, a final meal, a story told and kept: these honour what was attempted and help people carry it forward into whatever they build next.

\section{Genesis as a Ratified Workflow}

The Sovereign Stack can describe its own genesis. The phases in this chapter exist as a BPMN workflow template, kept in a shared commons of templates and playable in Oasis. A founding group can adapt it to its own circumstances, ratify it once its Layer 5 exists, and use it to track milestones, reminders, and decisions. In that sense the Stack helps to build the Stack.

This recursion is useful, but earlier drafts overstated it. The genesis template is not a self-assembling organism, and it is not a flawless or mathematically verifiable path to freedom. People perform every step. The Orchestrator tracks a template the community has chosen and ratified; it does not generate blueprints for people to sign. Every genesis, successful or not, teaches something that improves the template for the next group. What the protocol offers is not certainty but company: the accumulated experience of everyone who has tried before, offered to any group of people brave enough to begin.
